\chapter{Машина без дирижёра}
% полная версия: chapters/02_glutcomputer.tex

В вашем компьютере есть дирижёр --- тактовый генератор. Миллиарды раз в секунду он даёт отмашку, и все элементы схемы переключаются одновременно, в ногу. В мозге никакого дирижёра нет. Каждый нейрон работает сам по себе: сигналы приходят к нему когда придётся и уходят от него когда придётся. Давайте посмотрим, что можно вычислить в такой машине, если взять только \nt{ГЛУТ}-нейроны --- одни «плюсы».

\section*{Дырявое ведро}

Нейрон удобно представить ведром с дыркой. Каждый пришедший импульс --- это ковш воды. Вода понемногу вытекает через дырку. Если ковши приходят часто, вода успевает подняться до края --- и ведро переворачивается: нейрон щёлкает, ведро опустело, всё сначала. Если ковши приходят редко, вода успевает вытечь, и ничего не происходит.

\pencilfig{2}{\pencilart{2}}{Нейрон --- дырявое ведро: капли входов наполняют его, вода утекает; дошла до черты --- щелчок.}

Дырка --- главное отличие нейрона от логического вентиля. Нейрон помнит пришедший импульс не вечно, а какие-то миллисекунды. Всё, что пришло в пределах этого окна, складывается; то, что было раньше, забыто.

\section*{«Или» и «и»}

Если одного ковша хватает, чтобы перевернуть ведро, нейрон срабатывает на любой вход. Это «или». Если нужны два ковша, начинается интересное: «пришли ли оба?» без дирижёра не имеет смысла, пока не сказано --- \emph{в пределах какого времени}. Нейрон сработает, только если два импульса пришли почти одновременно, пока первый ковш не вытек. Это «и во времени», детектор совпадений. У одних нейронов окно --- несколько миллисекунд, у нейронов слуха --- доли миллисекунды.

\section*{Чего одни «плюсы» не умеют}

Попробуйте построить «не» --- нейрон, который работает, когда вход молчит. Не выйдет: тишина не наливает воду в ведро. И это верно не только для одного нейрона, но и для любой сети из одних «плюсов»: если входы усилить, ни один нейрон в такой сети не станет работать меньше. Отсюда самое неприятное свойство: \emph{активность нельзя погасить активностью}. Разгоревшийся пожар можно остановить только одним способом --- перестать подбрасывать дрова.

Природа нашла обходной путь. Во многих органах чувств один и тот же раздражитель передают две группы нейронов. В глазу одни клетки сообщают «стало светлее», другие --- «стало темнее». Если вместо одного провода у вас пара проводов «да» и «нет», то «не» получается даром: просто поменяйте провода местами. С такими парами из одних «плюсов» можно собрать любую логическую схему --- и даже понять, что вычисление закончено: в каждой паре загорелся ровно один провод. Инженеры строят так асинхронные схемы, которые правильно работают при любых задержках.

\section*{Время из задержек}

Без дирижёра нет общих часов, но время можно вычислять и так. Звук от источника сбоку приходит в ближнее ухо раньше, чем в дальнее, --- на доли миллисекунды. Проложим от левого уха провод слева направо вдоль ряда детекторов совпадений, а от правого --- справа налево. Два сигнала побегут навстречу друг другу и встретятся в той точке, куда оба добегут одновременно. Сработает детектор именно в этом месте. Разница во времени превратилась в \emph{номер} сработавшего нейрона --- то есть в направление на звук. Точность получается в десяток микросекунд, хотя каждый отдельный нейрон куда менее точен: выручает то, что детекторов много.

\pencilfig{102}{%
% delay line: the left-ear wire runs right along the top, the right-ear wire runs left along the bottom
\draw[pen=pcBlue] (0,0.6) -- (5.6,0.6); \draw[pen=pcBlue] (6.0,-0.6) -- (0.4,-0.6);
\foreach \i in {1,...,7}{
  \pgfmathsetmacro{\x}{0.75*\i}
  \draw[pen=pcBlue] (\x,0.6) -- (\x,0.24); \draw[pen=pcBlue] (\x,-0.6) -- (\x,-0.24);
  \ifnum\i=5 \draw[dotpen=pcAmber, wash=pcAmber, line width=1.1pt] (\x,0) circle (0.2);
  \else \draw[dotpen=pcBlue, wash=pcBlue] (\x,0) circle (0.2); \fi}
\draw[arr=pcBlue] (0.95,0.78) -- (1.75,0.78); \draw[arr=pcBlue] (5.3,-0.78) -- (4.5,-0.78);
\node[lbl, anchor=east] at (-0.05,0.6) {левое ухо};
\node[lbl, anchor=west] at (6.05,-0.6) {правое ухо};
\node[lbl, text=pcAmber!70!black, anchor=south] at (3.75,0.95) {здесь встретились};
\draw[arr=pcAmber!70!black] (3.75,0.95) -- (3.75,0.68);
% sound source on the left
\draw[penlite] (-1.25,-0.25) -- (-1.05,-0.25) -- (-0.8,-0.45) -- (-0.8,0.05) -- (-1.05,-0.15) -- (-1.25,-0.15) -- cycle;
\foreach \r in {0.2,0.35}{\draw[penlite] ($(-0.75,-0.2)+(-40:\r)$) arc (-40:40:\r);}
\node[lbl, anchor=north] at (-0.95,-0.5) {звук слева};
}{Звук слева раньше приходит в левое ухо, его сигнал успевает пробежать дальше, и встреча выходит справа от середины. Номер сработавшего нейрона --- направление на звук.}

\section*{Память --- это костёр}

Возьмём группу нейронов, которые сильно возбуждают друг друга. Если её разжечь коротким сигналом, она будет поддерживать сама себя, как костёр, в который огонь подбрасывает сам себя. Сигнал кончился, а группа всё горит --- она помнит, что сигнал был. Это ячейка памяти. А если сигнал был слишком коротким, костёр не успевает разгореться и гаснет.

Но вот беда: как стереть такую память? Активностью --- никак. Остаётся ждать, пока нейроны устанут и костёр прогорит.

\section*{Таймер, запущенный событием}

Соединим группы в цепочку: первая поджигает вторую, вторая третью. Толкнём первую --- и по цепочке побежит волна, как по костяшкам домино. Номер упавшей костяшки говорит, сколько прошло времени с толчка. Это не часы, которые тикают всегда, а таймер, который включается событием и отрабатывает один раз. У певчих птиц во время песни нейроны одного отдела срабатывают строго по очереди --- очень похоже на такую цепочку.

Итак, из одних «плюсов» получается многое. Не хватает трёх вещей: выбрать одно из многого, стереть память по команде и удержать сеть от пожара. Все три дают «минусы» --- о них следующая глава.

\fullversion{2}
